 Capítulo  5
 Interrupções  e  drivers  de  dispositivo
 Um  driver  é  o  código  em  um  sistema  operacional  que  gerencia  um  dispositivo  específico:  ele  configura  o  hardware  do  
dispositivo,  instrui  o  dispositivo  a  executar  operações,  lida  com  as  interrupções  resultantes  e  interage  com  processos  
que  podem  estar  aguardando  E/S  do  dispositivo.  O  código  do  driver  pode  ser  complicado,  pois  ele  é  executado  
simultaneamente  com  o  dispositivo  que  gerencia.  Além  disso,  o  driver  precisa  entender  a  interface  de  hardware  do  
dispositivo,  que  pode  ser  complexa  e  mal  documentada.
 Dispositivos  que  precisam  da  atenção  do  sistema  operacional  geralmente  podem  ser  configurados  para  gerar  
interrupções,  que  são  um  tipo  de  armadilha.  O  código  de  tratamento  de  armadilhas  do  kernel  reconhece  quando  
um  dispositivo  gera  uma  interrupção  e  chama  o  manipulador  de  interrupções  do  driver;  no  xv6,  esse  despacho  
ocorre  em  devintr  (kernel/trap.c:178).
 Muitos  drivers  de  dispositivo  executam  código  em  dois  contextos:  uma  metade  superior,  que  é  executada  na  thread  
do  kernel  de  um  processo,  e  uma  metade  inferior,  que  é  executada  em  tempo  de  interrupção.  A  metade  superior  é  
chamada  por  meio  de  chamadas  de  sistema,  como  leitura  e  gravação,  que  exigem  que  o  dispositivo  execute  E/S.  Esse  
código  pode  solicitar  ao  hardware  que  inicie  uma  operação  (por  exemplo,  solicitar  ao  disco  que  leia  um  bloco);  em  
seguida,  o  código  aguarda  a  conclusão  da  operação.  Por  fim,  o  dispositivo  conclui  a  operação  e  gera  uma  interrupção.  
O  manipulador  de  interrupção  do  driver,  atuando  como  a  metade  inferior,  descobre  qual  operação  foi  concluída,  ativa  
um  processo  em  espera,  se  apropriado,  e  informa  ao  hardware  para  iniciar  a  execução  de  qualquer  próxima  operação  em  espera.
 5.1  Código:  Entrada  do  console
 O  driver  do  console  (kernel/console.c)  é  uma  ilustração  simples  da  estrutura  do  driver.  O  driver  do  console  aceita  
caracteres  digitados  por  um  humano,  por  meio  do  hardware  de  porta  serial  UART  conectado  ao  RISC-V.
 O  driver  do  console  acumula  uma  linha  de  entrada  por  vez,  processando  caracteres  de  entrada  especiais,  como  
backspace  e  control-u.  Processos  do  usuário,  como  o  shell,  usam  a  chamada  de  sistema  read  para  buscar  linhas  de  
entrada  do  console.  Quando  você  digita  uma  entrada  para  o  xv6  no  QEMU,  suas  teclas  digitadas  são  enviadas  ao  xv6  
por  meio  do  hardware  UART  simulado  do  QEMU.
 O  hardware  UART  com  o  qual  o  driver  se  comunica  é  um  chip  16550  [13]  emulado  pelo  QEMU.  Em  um  computador  
real,  um  16550  gerenciaria  um  link  serial  RS232  conectado  a  um  terminal  ou  outro  computador.
 Ao  executar  o  QEMU,  ele  é  conectado  ao  seu  teclado  e  monitor.
 O  hardware  UART  parece,  para  o  software,  um  conjunto  de  registradores  de  controle  mapeados  na  memória.
 53
Machine Translated by Google
 Ou  seja,  existem  alguns  endereços  físicos  que  o  hardware  RISC-V  conecta  ao  dispositivo  UART,  de  modo  que  
cargas  e  armazenamentos  interagem  com  o  hardware  do  dispositivo  em  vez  da  RAM.  Os  endereços  mapeados  
na  memória  para  a  UART  começam  em  0x10000000,  ou  UART0  (kernel/memlayout.h:21).  Existem  vários  
registradores  de  controle  UART,  cada  um  com  a  largura  de  um  byte.  Seus  deslocamentos  em  relação  à  UART0  
são  definidos  em  (kernel/uart.c:22).  Por  exemplo,  o  registrador  LSR  contém  bits  que  indicam  se  os  caracteres  
de  entrada  estão  aguardando  para  serem  lidos  pelo  software.  Esses  caracteres  (se  houver)  estão  disponíveis  
para  leitura  no  registrador  RHR.  Cada  vez  que  um  caractere  é  lido,  o  hardware  da  UART  o  exclui  de  um  FIFO  
interno  de  caracteres  em  espera  e  limpa  o  bit  "pronto"  no  LSR  quando  o  FIFO  está  vazio.  O  hardware  de  
transmissão  da  UART  é  amplamente  independente  do  hardware  de  recepção;  se  o  software  escreve  um  byte  no  
THR,  a  UART  transmite  esse  byte.
 As  principais  chamadas  do  Xv6  consoleinit  (kernel/console.c:182)  para  inicializar  o  hardware  da  UART.  Este  
código  configura  a  UART  para  gerar  uma  interrupção  de  recebimento  quando  a  UART  recebe  cada  byte  de  
entrada  e  uma  interrupção  de  transmissão  completa  sempre  que  a  UART  termina  de  enviar  um  byte  de  saída  
(kernel/uart.c:53).
 O  shell  xv6  lê  o  console  por  meio  de  um  descritor  de  arquivo  aberto  por  init.c  (usuário/init.c:19).
 As  chamadas  para  o  sistema  read  passam  pelo  kernel  até  o  consoleread  (kernel/  console.c:80).  consoleread  
aguarda  a  chegada  da  entrada  (por  meio  de  interrupções)  e  seu  armazenamento  em  buffer  em  cons.buf,  copia  
a  entrada  para  o  espaço  do  usuário  e  (após  a  chegada  de  uma  linha  inteira)  retorna  ao  processo  do  usuário.  Se  
o  usuário  ainda  não  tiver  digitado  uma  linha  inteira,  qualquer  processo  de  leitura  aguardará  na  chamada  sleep  
(kernel/  console.c:96).  (O  Capítulo  7  explica  os  detalhes  do  sono).
 Quando  o  usuário  digita  um  caractere,  o  hardware  UART  solicita  ao  RISC-V  que  lance  uma  interrupção,  que  
ativa  o  manipulador  de  traps  do  xv6.  O  manipulador  de  traps  chama  devintr  (kernel/trap.c:178).  que  analisa  o  
registrador  RISC-V  Scause  para  descobrir  se  a  interrupção  é  de  um  dispositivo  externo.  Em  seguida,  solicita  a  
uma  unidade  de  hardware  chamada  PLIC  [3]  que  informe  qual  dispositivo  interrompeu  (kernel/trap.c:187).  Se  
fosse  o  UART,  devintr  chamaria  uartintr.
 uartintr  (kernel/uart.c:176)  lê  quaisquer  caracteres  de  entrada  em  espera  do  hardware  UART  e  os  entrega  
ao  consoleintr  (kernel/console.c:136);  Ele  não  espera  por  caracteres,  pois  entradas  futuras  gerarão  uma  nova  
interrupção.  A  função  do  consoleintr  é  acumular  caracteres  de  entrada  em  cons.buf  até  que  uma  linha  inteira  
chegue.  O  consoleintr  trata  o  backspace  e  alguns  outros  caracteres  de  forma  especial.  Quando  uma  nova  linha  
chega,  o  consoleintr  desperta  um  consoleread  em  espera  (se  houver).
 Uma  vez  ativado,  o  consoleread  observará  uma  linha  completa  em  cons.buf,  a  copiará  para  o  espaço  do  usuário  e  retornará  (por  
meio  do  mecanismo  de  chamada  do  sistema)  para  o  espaço  do  usuário.
 5.2  Código:  Saída  do  console
 Uma  chamada  de  sistema  de  gravação  em  um  descritor  de  arquivo  conectado  ao  console  eventualmente  
chega  em  uartputc  (kernel/uart.c:87).  O  driver  do  dispositivo  mantém  um  buffer  de  saída  (uart_tx_buf)  
para  que  os  processos  de  escrita  não  precisem  esperar  o  UART  terminar  de  enviar;  em  vez  disso,  o  
uartputc  anexa  cada  caractere  ao  buffer,  chama  o  uartstart  para  iniciar  a  transmissão  do  dispositivo  (se  
ainda  não  estiver)  e  retorna.  A  única  situação  em  que  o  uartputc  aguarda  é  se  o  buffer  já  estiver  cheio.
 Cada  vez  que  a  UART  termina  de  enviar  um  byte,  ela  gera  uma  interrupção.  uartintr  chama  uartstart,
 54
Machine Translated by Google
 que  verifica  se  o  dispositivo  realmente  concluiu  o  envio  e  entrega  ao  dispositivo  o  próximo  caractere  de  saída  armazenado  em  
buffer.  Assim,  se  um  processo  gravar  vários  bytes  no  console,  normalmente  o  primeiro  byte  será  enviado  pela  chamada  de  
uartputc  para  uartstart,  e  os  bytes  restantes  armazenados  em  buffer  serão  enviados  por  chamadas  de  uartstart  de  uartintr  à  
medida  que  as  interrupções  de  transmissão  completa  chegam.
 Um  padrão  geral  a  ser  observado  é  o  desacoplamento  da  atividade  do  dispositivo  da  atividade  do  processo  por  meio  de  
buffer  e  interrupções.  O  driver  do  console  pode  processar  a  entrada  mesmo  quando  nenhum  processo  está  aguardando  para  
lê-la;  uma  leitura  subsequente  verá  a  entrada.  Da  mesma  forma,  os  processos  podem  enviar  a  saída  sem  precisar  esperar  
pelo  dispositivo.  Esse  desacoplamento  pode  aumentar  o  desempenho,  permitindo  que  os  processos  sejam  executados  
simultaneamente  com  a  E/S  do  dispositivo,  e  é  particularmente  importante  quando  o  dispositivo  está  lento  (como  no  caso  da  
UART)  ou  precisa  de  atenção  imediata  (como  no  caso  de  caracteres  digitados  ecoando).  Essa  ideia  às  vezes  é  chamada  de  E/S.
 simultaneidade.
 5.3  Concorrência  em  drivers
 Você  deve  ter  notado  chamadas  para  acquire  em  consoleread  e  em  consoleintr.  Essas  chamadas  adquirem  um  bloqueio,  que  
protege  as  estruturas  de  dados  do  driver  de  console  contra  acesso  concorrente.  Há  três  perigos  de  simultaneidade  aqui:  dois  
processos  em  CPUs  diferentes  podem  chamar  consoleread  ao  mesmo  tempo;  o  hardware  pode  solicitar  que  uma  CPU  envie  
uma  interrupção  de  console  (na  verdade,  UART)  enquanto  essa  CPU  já  estiver  em  execução  dentro  de  consoleread;  e  o  
hardware  pode  enviar  uma  interrupção  de  console  em  uma  CPU  diferente  enquanto  consoleread  estiver  em  execução.  Esses  
perigos  podem  resultar  em  corridas  ou  deadlocks.  O  Capítulo  6  explora  esses  problemas  e  como  os  bloqueios  podem  resolvê
los.
 Outra  maneira  pela  qual  a  concorrência  requer  cuidado  nos  drivers  é  que  um  processo  pode  estar  aguardando  a  entrada  
de  um  dispositivo,  mas  a  chegada  da  sinalização  de  interrupção  da  entrada  pode  ocorrer  quando  um  processo  diferente  (ou  
nenhum  processo)  estiver  em  execução.  Assim,  os  manipuladores  de  interrupção  não  podem  pensar  sobre  o  processo  ou  
código  que  interromperam.  Por  exemplo,  um  manipulador  de  interrupção  não  pode  chamar  copyout  com  segurança  com  a  
tabela  de  páginas  do  processo  atual.  Os  manipuladores  de  interrupção  normalmente  realizam  relativamente  pouco  trabalho  
(por  exemplo,  apenas  copiam  os  dados  de  entrada  para  um  buffer)  e  ativam  o  código  da  metade  superior  para  fazer  o  resto.
 5.4  Interrupções  do  temporizador
 O  Xv6  utiliza  interrupções  de  temporizador  para  manter  seu  relógio  e  permitir  a  alternância  entre  processos  computacionais;  
as  chamadas  yield  em  usertrap  e  kerneltrap  causam  essa  alternância.  As  interrupções  de  temporizador  vêm  do  hardware  de  
relógio  conectado  a  cada  CPU  RISC-V.  O  Xv6  programa  esse  hardware  de  relógio  para  interromper  cada  CPU  periodicamente.
 O  RISC-V  exige  que  as  interrupções  de  timer  sejam  executadas  no  modo  máquina,  não  no  modo  supervisor.  O  modo  
máquina  RISC-V  executa  sem  paginação  e  com  um  conjunto  separado  de  registradores  de  controle,  portanto,  não  é  prático  
executar  código  kernel  xv6  comum  no  modo  máquina.  Como  resultado,  o  xv6  lida  com  interrupções  de  timer  completamente  
separadamente  do  mecanismo  de  captura  descrito  acima.
 O  código  executado  no  modo  de  máquina  em  start.c,  antes  do  main,  é  configurado  para  receber  interrupções  do  
temporizador  (kernel/start.c:63).  Parte  do  trabalho  consiste  em  programar  o  hardware  CLINT  (interruptor  local  do  núcleo)  para  
gerar  uma  interrupção  após  um  determinado  atraso.  Outra  parte  consiste  em  configurar  uma  área  de  risco,  análoga  à
 55
Machine Translated by Google
 trapframe,  para  ajudar  o  manipulador  de  interrupção  do  temporizador  a  salvar  registradores  e  o  endereço  dos  registradores  CLINT.
 Por  fim,  start  define  mtvec  como  timervec  e  habilita  interrupções  de  timer.
 Uma  interrupção  de  temporizador  pode  ocorrer  a  qualquer  momento  durante  a  execução  de  código  do  usuário  ou  do  kernel;  
não  há  como  o  kernel  desabilitar  interrupções  de  temporizador  durante  operações  críticas.  Portanto,  o  manipulador  de  interrupções  
de  temporizador  deve  executar  sua  função  de  forma  a  garantir  que  não  perturbe  o  código  do  kernel  interrompido.  A  estratégia  básica  
é  que  o  manipulador  solicite  ao  RISC-V  que  gere  uma  "interrupção  de  software"  e  retorne  imediatamente.  O  RISC-V  entrega  
interrupções  de  software  ao  kernel  com  o  mecanismo  de  captura  comum  e  permite  que  o  kernel  as  desabilite.  O  código  para  lidar  
com  a  interrupção  de  software  gerada  por  uma  interrupção  de  temporizador  pode  ser  visto  em  devintr  (kernel/trap.c:205).
 O  manipulador  de  interrupção  do  temporizador  no  modo  máquina  é  timervec  (kernel/kernelvec.S:95).  Ele  salva  alguns  
registradores  na  área  de  rascunho  preparada  por  start,  informa  ao  CLINT  quando  gerar  a  próxima  interrupção  do  
temporizador,  solicita  ao  RISC-V  que  lance  uma  interrupção  de  software,  restaura  os  registradores  e  retorna.  Não  há  
código  C  no  manipulador  de  interrupção  do  temporizador.
 5.5  Mundo  real
 O  Xv6  permite  interrupções  de  dispositivo  e  temporizador  durante  a  execução  no  kernel,  bem  como  durante  a  execução  de  
programas  do  usuário.  As  interrupções  de  temporizador  forçam  uma  troca  de  thread  (uma  chamada  para  yield)  a  partir  do  
manipulador  de  interrupções  de  temporizador,  mesmo  durante  a  execução  no  kernel.  A  capacidade  de  dividir  o  tempo  da  CPU  de  
forma  justa  entre  as  threads  do  kernel  é  útil  se  as  threads  do  kernel  às  vezes  gastam  muito  tempo  computando,  sem  retornar  ao  
espaço  do  usuário.  No  entanto,  a  necessidade  de  o  código  do  kernel  estar  ciente  de  que  pode  ser  suspenso  (devido  a  uma  
interrupção  de  temporizador)  e  posteriormente  retomado  em  uma  CPU  diferente  é  a  fonte  de  alguma  complexidade  no  xv6  (consulte  
a  Seção  6.6).  O  kernel  poderia  ser  um  pouco  mais  simples  se  as  interrupções  de  dispositivo  e  temporizador  ocorressem  apenas  
durante  a  execução  do  código  do  usuário.
 Oferecer  suporte  a  todos  os  dispositivos  em  um  computador  típico  em  todo  o  seu  esplendor  dá  muito  trabalho,  pois  
há  muitos  dispositivos,  os  dispositivos  têm  muitos  recursos  e  o  protocolo  entre  o  dispositivo  e  o  driver  pode  ser  complexo  
e  mal  documentado.  Em  muitos  sistemas  operacionais,  os  drivers  ocupam  mais  código  do  que  o  kernel  principal.
 O  driver  UART  recupera  dados  um  byte  por  vez,  lendo  os  registradores  de  controle  UART;  esse  padrão  é  chamado  de  E/S  
programada,  pois  o  software  controla  a  movimentação  dos  dados.  A  E/S  programada  é  simples,  mas  lenta  demais  para  ser  usada  
em  altas  taxas  de  dados.  Dispositivos  que  precisam  mover  grandes  quantidades  de  dados  em  alta  velocidade  normalmente  usam  
acesso  direto  à  memória  (DMA).  O  hardware  do  dispositivo  DMA  grava  os  dados  de  entrada  diretamente  na  RAM  e  lê  os  dados  de  
saída  da  RAM.  Dispositivos  de  disco  e  rede  modernos  usam  DMA.
 Um  driver  para  um  dispositivo  DMA  prepararia  dados  na  RAM  e,  então,  usaria  uma  única  gravação  em  um  registro  de  controle  para  
dizer  ao  dispositivo  para  processar  os  dados  preparados.
 Interrupções  fazem  sentido  quando  um  dispositivo  precisa  de  atenção  em  momentos  imprevisíveis  e  não  com  muita  frequência.
 Mas  as  interrupções  têm  alta  sobrecarga  de  CPU.  Assim,  dispositivos  de  alta  velocidade,  como  controladores  de  rede  e  de  disco,  
usam  truques  que  reduzem  a  necessidade  de  interrupções.  Um  truque  é  gerar  uma  única  interrupção  para  um  lote  inteiro  de  
solicitações  de  entrada  ou  saída.  Outro  truque  é  o  driver  desabilitar  completamente  as  interrupções  e  verificar  o  dispositivo  
periodicamente  para  ver  se  ele  precisa  de  atenção.  Essa  técnica  é  chamada  de  polling.  O  polling  faz  sentido  se  o  dispositivo  executa  
operações  muito  rapidamente,  mas  desperdiça  tempo  de  CPU  se  o  dispositivo  estiver  ocioso  na  maior  parte  do  tempo.  Alguns  
drivers  alternam  dinamicamente  entre  polling  e  interrupções.
 56
Machine Translated by Google
 dependendo  da  carga  atual  do  dispositivo.
 O  driver  UART  copia  os  dados  recebidos  primeiro  para  um  buffer  no  kernel  e  depois  para  o  espaço  do  usuário.
 Isso  faz  sentido  em  baixas  taxas  de  dados,  mas  essa  cópia  dupla  pode  reduzir  significativamente  o  desempenho  de  
dispositivos  que  geram  ou  consomem  dados  muito  rapidamente.  Alguns  sistemas  operacionais  conseguem  mover  dados  
diretamente  entre  os  buffers  do  espaço  do  usuário  e  o  hardware  do  dispositivo,  geralmente  com  DMA.
 Conforme  mencionado  no  Capítulo  1,  o  console  aparece  para  os  aplicativos  como  um  arquivo  comum,  e  os  aplicativos  
leem  a  entrada  e  escrevem  a  saída  usando  as  chamadas  de  sistema  de  leitura  e  escrita .  Os  aplicativos  podem  querer  
controlar  aspectos  de  um  dispositivo  que  não  podem  ser  expressos  pelas  chamadas  de  sistema  de  arquivos  padrão  (por  
exemplo,  habilitar/desabilitar  o  buffer  de  linha  no  driver  do  console).  Os  sistemas  operacionais  Unix  suportam  a  chamada  de  
sistema  ioctl  para  esses  casos.
 Alguns  usos  de  computadores  exigem  que  o  sistema  responda  em  um  tempo  limitado.  Por  exemplo,  em  sistemas  
críticos  de  segurança,  perder  um  prazo  pode  levar  a  desastres.  O  Xv6  não  é  adequado  para  configurações  de  tempo  real  
rígido.  Sistemas  operacionais  para  tempo  real  rígido  tendem  a  usar  bibliotecas  que  se  conectam  ao  aplicativo  de  forma  a  
permitir  uma  análise  para  determinar  o  pior  tempo  de  resposta.  O  Xv6  também  não  é  adequado  para  aplicativos  de  tempo  
real  flexível,  quando  perder  um  prazo  ocasionalmente  é  aceitável,  porque  o  escalonador  do  Xv6  é  muito  simplista  e  possui  
um  caminho  de  código  do  kernel  onde  as  interrupções  são  desabilitadas  por  um  longo  período.
